\documentclass[a4paper,10pt]{article}


\usepackage[utf8]{inputenc}
\usepackage[italian]{babel}
\usepackage{amsfonts}
\usepackage{amsmath}
\usepackage{amssymb}
\usepackage{amsthm}
\usepackage{fancyhdr}
\usepackage{fullpage}
\usepackage{graphicx}
\usepackage[colorlinks]{hyperref}
\usepackage{makeidx}
\usepackage{mathpazo}
\usepackage{multicol}
\usepackage[parfill]{parskip}
%\usepackage{utopia}

% * define a `\twoidxcolumn` based on `\twocolumn`:
\def\twoidxcolumn{%
%\clearpage
\global\columnwidth\textwidth%
\global\advance\columnwidth-\columnsep%
\global\divide\columnwidth\tw@%
\global\hsize\columnwidth%
\global\linewidth\columnwidth%
\global\@twocolumntrue%
\global\@firstcolumntrue%
\col@number \tw@%
%\@ifnextchar [\@topnewpage
\@floatplacement%
}%


\newcommand{\smallspace}{\vskip0.3cm}

% Title Page
\title{Lupus in tempo reale \\ Regolamento della dodicesima edizione}
\author{Alessandro Iraci}

\begin{document}

\maketitle

\begin{tabular}{lp{0.6\textwidth}}
      \begin{minipage}{0.22\textwidth}
            \vspace{3mm}
            \href{http://creativecommons.org/licenses/by/4.0/}{\includegraphics{ccby.pdf}}
      \end{minipage}
       &
      This work is licensed under a \href{http://creativecommons.org/licenses/by/4.0/}{Creative Commons Attribution 4.0 International License}.
\end{tabular}


\section{Introduzione}

\subsection{Cos'è Lupus in tempo reale?}

\emph{Lupus in tempo reale} è una variante dei più tradizionali \emph{Lupus in Tabula} e \emph{Mafia}, caratterizzata da un maggiore coinvolgimento e da interazioni più sofisticate tra i giocatori.
La principale differenza consiste nel ritmo del gioco: la partita si estende infatti su diversi giorni, in modo aderente all'ambientazione.
Sono premiate dalla dinamica di gioco la logica, la capacità di bluffare (e di smascherare i bluff altrui), la capacità di saper coordinare un'azione di gruppo, il carisma e le velleità leaderistiche.

\emph{Lupus in tempo reale} è stato inventato nel 2010, ed è stato giocato una volta all'anno fino al 2012 sotto la supervisione di Francesco Guatieri come principale Game Master. A partire dall'anno 2014 è passato sotto il controllo di Giovanni Mascellani e Giovanni Paolini, e da lì in poi l'organigramma si è esteso edizione per edizione (vedi Sezione~\ref{credits}).

Nel seguito di questo regolamento, verrà utilizzata la sigla ``GM'' come abbreviazione di ``Game Master''.

\subsection{Ambientazione}

La storia si svolge in quello che una volta era un tranquillo villaggio di campagna. A turbare la quiete del posto è un evento quantomeno atipico, ossia un'invasione di lupi mannari. Nascosti sotto le spoglie di normalissimi contadini, i licantropi hanno intenzione di prendere il controllo del villaggio e sbranare chiunque non sia disposto ad allearsi con loro.

Le autorità nazionali, venendo a sapere dell'accaduto e temendo il diffondersi della licantropia in tutto il paese, hanno disposto un blocco forzato delle vie d'accesso al villaggio, rendendo virtualmente impossibile a chiunque entrare o uscire dalla frazione.

Per tentare di ristabilire la normalità viene imposta una durissima legge marziale: ogni sera, il villaggio convocherà un'assemblea composta da tutti gli abitanti, al termine della quale un sospetto lupo mannaro verrà arso sul rogo. La sfiducia nei confronti degli altri abitanti è totale, chiunque potrebbe essere un licantropo sotto mentite spoglie. La propria vicina di casa, l'insospettabile anziano signore dell'ortofrutta, o anche quel nuovo arrivato peloso e con le orecchie a punta, così normali di giorno, potrebbero col calar della notte assumere le sembianze di un feroce lupo assetato di sangue.

Fra i cittadini c'è paura, e tanta indecisione su quale sia la cosa giusta da fare. Meglio combattere i lupi mannari e i loro alleati e cercare di liberare il villaggio, oppure allearsi con loro nella speranza di avere salva la vita, e magari anche di ottenere una posizione di potere? Meglio rischiare la pelle andando ad origliare dietro la porta del proprio vicino, cercando di capire se questi è un lupo, oppure stare tranquilli nel proprio letto sperando di non ricevere visite sgradite? Molti abitanti hanno deciso di attivarsi per riprendere il controllo del villaggio; altri hanno ceduto, per paura della morte o per sete di potere, alle lusinghe dei licantropi, aiutandoli a portare a termine i loro progetti.

Ciascuno ha tutto l'interesse a tenere nascosta la propria scelta: per i traditori, venire scoperti potrebbe significare morte sul rogo; per i villici che hanno deciso di combattere l'invasione, esporsi troppo potrebbe voler dire essere la prossima vittima dei lupi mannari, oppure ricevere terrificanti visite da parte di uno spettro. Ma non sempre tacere per la propria incolumità è la scelta giusta. In determinate circostanze rivelare preziose informazioni, anche a costo di rischiare la vita, potrebbe essere fondamentale: veder trionfare i lupi mannari significa morte in ogni caso.

Qualunque sia stata la propria decisione, tuttavia, la maggior parte dei cittadini si attiva durante la notte per raccogliere informazioni o infastidire i nemici: per esempio appostandosi fuori da una casa per cercare di capire cosa succede all'interno, utilizzando le proprie capacità divinatorie per scoprire l'identità degli altri, girando armati per proteggere i concittadini dall'attacco dei licantropi, o compiendo rituali per tenere lontani gli spettri. Ciascuno fa il possibile, ma anche gli alleati di lupi mannari si mobilitano per intralciare le indagini, diffondere false informazioni e manipolare le votazioni durante l'assemblea.

Nel villaggio, ormai, la vita è una questione di doppio gioco; ogni passo falso rischia di costare la pelle non solamente a sé stessi, ma anche a tutti i propri alleati. Un delicatissimo equilibrio coinvolge le due parti, cercare di spezzarlo in proprio favore è l'obiettivo di tutti. Solo i più abili manipolatori riusciranno a prendere in mano la situazione e condurre la propria squadra alla vittoria. Lo scopo da raggiungere è chiaro, ma quale sia la fazione che riuscirà a trionfare dipende da voi!

\subsection{Partecipanti}

Possono partecipare:

\begin{itemize}
      \item gli studenti della Scuola;
      \item altre persone che frequentano spesso gli ambienti della Scuola, in particolare il collegio Faedo (solitamente) e la mensa.
\end{itemize}

Tutti i partecipanti si impegnano a farsi vedere un po' in giro negli ambienti della Scuola nei giorni della partita. Nessuno vi chiede di non partire venerdì mattina per tornare a casa se non avete lezione, ma se vivete da reclusi in camera vostra e non mangiate mai in mensa potrebbe essere il caso di non giocare.

% 
% \subsection{Quando si giocherà}
% 
% La partita inizierà la sera di martedì 18 novembre 2014, con ritrovo al collegio Carducci. % TODO

\clearpage

% \subsection{Principali cambiamenti nella quarta edizione}
% 
% Questa sezione è pensata per dare a chi ha giocato alla terza edizione di Lupus
% un'idea delle modifiche più importanti che sono state apportate. Tuttavia non
% sostituisce un'attenta lettura del resto del regolamento.
% 
% \begin{itemize}
%  \item La Fazione dei Vampiri è stata sostituita dalla nuova Fazione dei
% Negromanti, che ha meccaniche di gioco molto diverse.
%  \item È stato eliminato il suicidio.
%  \item È stato eliminato il rituale dei fantasmi.
%  \item I possibili colori dell'aura sono tornati ad essere solamente il bianco e
% il nero.
%  \item Molti ruoli sono stati creati, eliminati o modificati. A chi è abituato
% al regolamento della terza edizione raccomandiamo di controllare con particolare
% attenzione i seguenti ruoli, che hanno subito i cambiamenti più rilevanti:
% Cacciatore, Custode del cimitero, Esorcista, Investigatore, Rinnegato, Profanatore di tombe, Negromante, Fantasma, Ipnotista, Medium, Spettro.
%  \item Le comunicazioni di gioco (votazione, utilizzo dei poteri, resoconto
% dell'alba e del tramonto) avvengono attraverso un'interfaccia web e non più via
% e-mail.
% \end{itemize}
% 
% \subsection{Principali cambiamenti nella quinta edizione}
% 
% Le seguenti sono le modifiche apportate rispetto alla quarta edizione di Lupus in tempo reale.
% 
% \begin{itemize}
% \item È stato introdotto lo Scrutatore.
% \item È stato eliminato lo Spettro con il potere della Duplicazione.
% \item L'ultimo Ipnotista, quando muore, diventa uno Spettro con il potere dell'Ipnosi.
% \item Il Profanatore di Tombe è stato trasferito alla Fazione dei
%   Popolani ed è stato rinominato in Sciamano. 
%   Rimane con aura nera, ma diventa mistico.
% \item Il Necrofilo è stato rinominato in Trasformista.
% \item Il Medium scopre il ruolo del personaggio scelto, e non più l'aura.
% \end{itemize}

% 
% 
% \subsection{Principali cambiamenti nella sesta edizione}
% 
% Questa sezione è pensata per dare a chi ha giocato alla quinta edizione di Lupus un'idea delle modifiche più importanti che sono state apportate. Tuttavia non sostituisce un'attenta lettura del resto del regolamento.
% 
% \begin{itemize}
%   \item Il Trasformista e il Fantasma hanno ora aura nera.
%   \item Nessun personaggio può utilizzare il proprio potere su sé stesso.
%   \item Sono state rimosse le restrizioni sull'utilizzo di un potere due notti consecutive sullo stesso personaggio.
%   \item Sono stati eliminati il Custode del Cimitero, il Cacciatore, e il potere soprannaturale della Mistificazione.
%   \item Sono stati introdotti l'Assassino, lo Stregone, e il potere soprannaturale della Confusione.
%   \item I personaggi della Fazione dei Lupi non possono diventare Spettri.
%   \item Le conoscenze iniziali dei personaggi della Fazione dei Lupi sono state modificate.
%   \item Il potere soprannaturale dell'Ipnosi può ora essere scelto dai Negromanti; il primo Ipnotista che muore diventa uno Spettro con tale potere (se disponibile).
%   \item Il Messia fallisce se utilizza il proprio potere su uno Spettro.
%   \item L'Avvocato del Diavolo è stato rinominato in Avvocato.
% \end{itemize}
%
%\subsection{Principali cambiamenti nella settima edizione}
% 
%Questa sezione è pensata per dare a chi ha giocato alla sesta edizione di Lupus un'idea delle modifiche più importanti che sono state apportate. Tuttavia non sostituisce un'attenta lettura del resto del regolamento.
% 
%\begin{itemize}
%	\item Il poteri speciali del Trasformista, dei Lupi, ed il potere soprannaturale della Morte hanno effetto solo sui componenti della Fazione dei Popolani.
% 	\item Il potere speciale del Diavolo non ha effetto sui componenti della Fazione dei Negromanti.
% 	\item Il potere speciale dei Negromanti può essere usato ogni notte.
% 	\item L'Ipnotista non diviene uno Spettro dopo la sua morte.
% 	\item I seguenti poteri soprannaturali sono stati modificati: Confusione, Ipnosi, Morte, Visione. È stato aggiunto un nuovo potere soprannaturale: Corruzione.
% 	\item Le restrizioni sulla spettrificazione sono state modificate.
% 	\item È stato reintrodotto il ruolo di Custode del cimitero. Oltre al vecchio potere, Custode del cimitero e Guardia del corpo scoprono anche quanti personaggi hanno agito sul loro stesso bersaglio.
% 	\item È stato introdotto un nuovo ruolo: il Necrofilo. Appartiene alla Fazione dei Lupi.
% 	\item La Fattucchiera ha aura bianca.
% 	\item Il potere speciale dello Scrutatore è stato modificato.
%\end{itemize}

%     \subsection{Principali cambiamenti nell'ottava edizione}
% 
%     Questa sezione è pensata per dare a chi ha giocato alla settima edizione di Lupus un'idea delle modifiche più importanti che sono state apportate. Tuttavia non sostituisce un'attenta lettura del resto del regolamento.
% 
%     \begin{itemize}
% 	    \item Le meccaniche di spettrificazione sono state modificate, così come il funzionamento degli Spettri.
% 	    \item Gli Spettri sono ora associati ad Incantesimi e non a poteri soprannaturali.
% 	    \item I poteri speciali si chiamano ora \emph{abilità}. I poteri soprannaturali si chiamano ora semplicemente \emph{poteri}.
% 	    \item I fallimenti non contano più come utilizzo di abilità o potere se il suo utilizzo è permesso ogni due notti, oppure una sola volta a partita.
% 	    \item Per uccidere qualcuno sul rogo, occorre che questi riceva più del 50\% dei voti disponibili. Inoltre, i voti espressi per il rogo adesso sono segreti; il totale dei voti ricevuti da ciascuno rimane pubblico.
% 	    \item La carica di Sindaco è stata eliminata.
% 	    \item L'alba è stata spostata alle ore 9:00.
% 	    \item Tutti i componenti della Fazione dei Lupi si conoscono fra di loro. Lo stesso vale per la Fazione dei Negromanti.
% 	    \item Sono state eliminate le restrizioni sull'uccisione dei componenti delle Fazioni dei Lupi e dei Negromanti relative all'altra Fazione.
% 	    \item Sono stati eliminati i seguenti ruoli: Custode del cimitero, Avvocato, Necrofilo, Rinnegato, Ipnotista, Medium, Scrutatore.
% 	    \item Sono stati modificati i seguenti ruoli: Divinatore, Esorcista, Investigatore, Sciamano, Trasformista, Diavolo, Fattucchiera, Sequestratore, Negromante.
% 	    \item Sono stati introdotti i seguenti ruoli: Cacciatore, Spia, Alcolista.
% 	    \item Sono stati eliminati i seguenti Incantesimi: Ipnosi, Visione.
% 	    \item Sono stati modificati i seguenti Incantesimi: Amnesia, Confusione, Morte, Occultamento.
% 	    \item Sono stati introdotti i seguenti Incantesimi: Assoluzione, Diffamazione, Telepatia, Vita.
%     \end{itemize}

% \subsection{Principali cambiamenti nella decima edizione}

% Questa sezione è pensata per dare a chi ha giocato alla ottava o alla nona edizione di Lupus un'idea delle modifiche più importanti che sono state apportate. Tuttavia non sostituisce un'attenta lettura del resto del regolamento.

% \begin{itemize}
% 	\item La Fazione dei Negromanti è stata eliminata.
% 	\item I ruoli di Negromante e Fantasma fanno ora parte della Fazione dei Lupi.
% 	\item È stata introdotta la possibilità che alcuni giocatori ricevano informazioni iniziali aggiuntive oltre a quelle pubbliche.
% 	\item Sono stati eliminati i seguenti ruoli: Spia, Massone, Alcolista.
% 	\item Sono stati modificati i seguenti ruoli: Esorcista, Sciamano, Stregone, Negromante, Fantasma.
% 	\item Sono stati introdotti i seguenti ruoli: Apprendista.
% 	\item Gli Incantesimi sono stati modificati per riadattarli alla fazione dei Lupi.
% \end{itemize}

% \subsection{Principali cambiamenti nell'undicesima edizione}

% Questa sezione è pensata per dare a chi ha giocato alla decima edizione di Lupus un'idea delle modifiche più importanti che sono state apportate. Tuttavia non sostituisce un'attenta lettura del resto del regolamento.

% \begin{itemize}
%     \item La Fazione dei Popolani si chiama ora Fazione dei Villici.
%     \item La Fazione dei Lupi si chiama ora Fazione dei Villani.
%     \item Sono stati eliminati i seguenti ruoli: Investigatore, Sciamano.
%     \item Sono stati modificati i seguenti ruoli: Apprendista, Cacciatore, Esorcista, Espansivo, Guardia, Trasformista, Assassino, Fantasma, Fattucchiera, Negromante, Sequestratore, Stregone.
%     \item Sono stati introdotti i seguenti ruoli: Alcolista, Medium.
%     \item È stata modificata la dinamica di assegnazione degli Incantesimi, nonché diversi Incantesimi.
%     \item Il ``giorno lungo'' è stato sostituito da due ``notti lunghe'', mercoledì e sabato.
% \end{itemize}

\subsection{Principali cambiamenti nella dodicesima edizione}

Questa sezione è pensata per dare a chi ha giocato all'undicesima edizione di Lupus un'idea delle modifiche più importanti che sono state apportate. Tuttavia non sostituisce un'attenta lettura del resto del regolamento.

\begin{itemize}
      \item È stato modificato il funzionamento del rogo.
      \item La Fazione dei Lupi si chiama ora Fazione dei Villani.
      \item Sono stati eliminati i seguenti ruoli: Investigatore, Sciamano.
      \item Sono stati modificati i seguenti ruoli: Apprendista, Cacciatore, Esorcista, Espansivo, Guardia, Trasformista, Assassino, Fantasma, Fattucchiera, Negromante, Sequestratore, Stregone.
      \item Sono stati introdotti i seguenti ruoli: Alcolista, Medium.
      \item È stata modificata la dinamica di assegnazione degli Incantesimi, nonché diversi Incantesimi.
      \item Il ``giorno lungo'' è stato sostituito da due ``notti lunghe'', mercoledì e sabato.
\end{itemize}

\pagebreak

\section{Metagioco}

\subsection{Giorni e notti di gioco}

I giorni di gioco si estendono (salvo eccezioni concordate a inizio partita) dalle 9:00 alle 22:00 (circa) di lunedì, martedì, giovedì, venerdì e domenica alle 22:00 della sera stessa. Le notti di gioco si estendono dalle 22:00 (circa) dei giorni indicati fino all'inizio del giorno di gioco successivo. Il mercoledì e il sabato sono giorni di pausa, di cui i giocatori possono approfittare per dedicarsi ad altro.

\subsection{Interfaccia web}

Prima dell'inizio della partita, ciascun giocatore deve crearsi un account per accedere all'interfaccia web di Lupus, il cui indirizzo è \verb|http://lupus.uz.sns.it/|.
Le votazioni durante il giorno e le attivazioni di abilità e poteri avvengono tutte tramite l'interfaccia web. Le informazioni sulle votazioni del giorno e sugli avvenimenti della notte compaiono a loro volta sull'interfaccia web (le informazioni pubbliche possono essere visualizzate senza bisogno di autenticazione, mentre quelle private sono accessibili solo dopo il login). In caso di impossibilità di accedere a internet, i giocatori possono eccezionalmente contattare i GM per comunicare loro le proprie intenzioni di voto o il modo in cui desiderano utilizzare le proprie abilità o poteri.

\subsection{Restrizioni}

Nel rispetto dello spirito del gioco, è opportuno introdurre alcune restrizioni.

Le comunicazioni fra GM e giocatori (siano esse verbali, scritte, digitali, telepatiche, o in qualsiasi altra forma) sono inviolabili. Lo stesso vale per l'utilizzo dell'area privata dell'interfaccia web, le lettere cartacee consegnate all'inizio della partita, l'account email ufficiale, e l'account del forum che essi designano come ufficiale.
Con \emph{inviolabili} si intende che tali mezzi non possono essere letti, mostrati, registrati, né utilizzati in qualsiasi modo durante discussioni pubbliche o private con parti terze o in modi che parimenti violino lo spirito del gioco.

Gli altri mezzi di comunicazione dei giocatori, quali i loro cellulari, personal computer, o account di messaggistica di sorta, sono considerati privati.
Con \emph{privati} si intende che i giocatori non possono essere incoraggiati, minacciati, o costretti in qualsiasi modo a mostrare o concedere l'accesso a tali mezzi a parti terze, o utilizzarli in modi che parimenti violino lo spirito del gioco. Gli account anonimi del forum sono parimenti considerati privati, e le restrizioni sopra citate si applicano anche al dimostrare la proprietà di un account o mostrarne i contenuti.

In sintesi, è vietato cercare di ottenere informazioni su di un giocatore violando o cercando di violare uno qualsiasi dei mezzi citati, e più in generale non è consentito effettuare metagioco sulle comunicazioni private dei giocatori.

L'utilizzo di qualsiasi forma di crittografia non è consentito. Sono parimenti vietate altre forme di controllo o registrazione delle azioni dei giocatori in modi che violino lo spirito del gioco.

Ci possono essere altri comportamenti che non ricadono in quelli descritti in precedenza ma che comunque violano lo spirito del gioco. In caso di dubbi, chiedete preventivamente ai GM.

In ciascuno di questi casi i GM hanno il diritto di squalificare un giocatore o un'intera fazione, e il loro giudizio è inappellabile.

Infine, per mitigare l'effetto di giocatori inattivi, se un giocatore non accede all'interfaccia di gioco per quattro sessioni consecutive (due giorni e due notti), salvo giustificazione richiesta ai GM, tale giocatore viene immediatamente squalificato. I GM possono concedere, a loro discrezione, dei piccoli bonus alla fazione di appartenenza di tale giocatore per mantenere l'equilibrio nella partita.

\subsection{Comunicazioni tra i giocatori}

I giocatori possono comunicare tra di loro con qualsiasi mezzo. Possono dirsi qualsiasi cosa.

A disposizione dei giocatori viene messo un forum, raggiungibile tramite l'interfaccia web della partita.  All'inizio della partita, verrà pubblicata una lista degli account ufficiali dei partecipanti. Ai giocatori è vietato usare un account di quella lista diverso dal proprio: in particolare, è vietato mandare messaggi da quell'account se vi si ha accesso in qualche modo; è vietato usare qualsiasi tipo di abilità informatica per avere accesso a quell'account del forum; è vietato inviare messaggi usando quell'account pur non avendone accesso (ma se sapete come fare, vi preghiamo di segnalarlo così che possiamo impedire che possa accadere).

La violazione di queste regole può portare alla squalifica del giocatore o, nei casi più gravi, dell'intera fazione. La violazione delle leggi dello Stato in cui vi trovate è sconsigliata dai GM, ma non è proibita a priori dal regolamento. Questo non vi darà comunque alcun tipo di immunità, civile o penale.

È lecito (anzi, spesso consigliato) imbrogliare gli altri giocatori, mentire, usare account del forum fasulli o anonimi per comunicare con gli altri giocatori o per cercare di ingannarli (in particolare, c'è una sezione del forum appositamente dedicata), e qualsiasi altro mezzo vi venga in mente per vincere la partita, purché sia nel rispetto delle regole di cui sopra.

È anche lecito spiare il comportamento di altri giocatori, sia di persona che con mezzi informatici, con la seguente eccezione: è vietato spiare un giocatore all'interno della sua camera o abitazione personale. In particolare è vietato eseguire monitoraggio delle reti dei collegi, allo scopo per esempio di determinare in che istanti i giocatori sono connessi ad Internet.

Anche i morti possono parlare con gli altri membri del villaggio, e continuano a giocare. Si incoraggiano i giocatori a non comunicare solo tramite il forum ma anche e soprattutto di persona.

Naturalmente le regole sopra indicate si possono prestare a diverse interpretazioni. I giocatori sono invitati a chiedere ai GM conferma della liceità di comportamenti che potrebbero essere valutati \emph{borderline}.

\pagebreak

\section{Gioco}

\subsection{Fazioni e condizioni di vittoria}

I giocatori sono divisi in due fazioni: la Fazione dei Villici e la Fazione dei Villani. I Villani si conoscono fra di loro sin dall'inizio della partita.

Una fazione vince se, subito dopo l'alba oppure subito dopo la votazione del tramonto, tutti i personaggi vivi appartengono a quella fazione. Inoltre, i Villici vincono immediatamente se muoiono tutti i Lupi.

Se una fazione ha chiaramente vinto prima che le condizioni di vittoria precedenti siano rispettate, i GM possono porre fine alla partita in anticipo (ma non devono farlo per forza).

\subsection{Inizio del gioco}

All'inizio della partita ciascun giocatore riceve le informazioni necessarie per partecipare, tra cui il ruolo ed eventualmente le identità di alcuni altri giocatori. La partita comincia con la notte.


\subsection{Svolgimento del giorno e votazioni}

L'inizio del giorno è annunciato sul sito web, insieme a tutte le informazioni ottenute da ciascun giocatore in seguito agli avvenimenti della notte appena trascorsa. Ogni giorno, il villaggio convoca un'assemblea per stabilire chi condannare a morte sul rogo. I personaggi vivi hanno il diritto di votare entro le ore 22:30 un abitante del villaggio da condannare al rogo, oppure di votare perché non ci sia alcun rogo. Dopo le ore 22:30, viene pubblicata la lista delle persone votate, insieme al numero di preferenze ricevute da ciascun personaggio.

Se un personaggio riceve strettamente più del 50\% dei voti espressi, purché strettamente più del 25\% del numero dei giocatori vivi, viene condannato al rogo e muore. Altrimenti, nessuno viene condannato.

\subsection{Svolgimento della notte e uso di abilità e poteri}

La notte inizia nel momento in cui viene pubblicato l'esito della votazione del giorno, e termina alle 9:00 della mattina seguente. Con \emph{abilità} si intende la capacità di compiere un'azione speciale da parte di un personaggio vivo; con \emph{potere} si intende la stessa cosa, ma da parte di un personaggio morto. Durante la notte, i personaggi vivi che hanno delle abilità hanno la facoltà di attivarle, ed i personaggi morti che hanno dei poteri possono fare lo stesso.

Abilità e poteri possono essere usati con frequenza diversa, per compensare la diversa forza.
\begin{itemize}
      \item Le abilità e i poteri utilizzabili ``ogni notte'' possono essere usati senza restrizioni.
      \item Le abilità e i poteri utilizzabili ``ogni due notti'', quando utilizzati con successo, impediscono di utilizzare altre abilità o poteri nella notte successiva.
      \item Le abilità e i poteri utilizzabili ``una volta a partita'', quando utilizzati con successo, impediscono di utilizzare altre abilità o poteri per il resto della partita.
\end{itemize}

Le abilità con il tag [morte] non possono essere utilizzate durante la prima notte di gioco. Nessun personaggio può utilizzare la propria abilità o potere su sé stesso.


%Un personaggio che usa la propria abilità o il proprio potere su un altro personaggio, 

\subsection{Movimenti}

Durante la notte, se un personaggio (vivo) utilizza la propria abilità su un altro personaggio, genera un proprio movimento verso la casa del personaggio (vivo o morto) su cui tale abilità viene utilizzata. L'uso di un potere (da parte di un personaggio morto), invece, non genera alcun movimento. Alcuni ruoli possono ottenere informazioni sui movimenti o cancellarli.

Come specificato nella sezione \ref{fallimento}, il fallimento nell'utilizzo di un'abilità genera comunque un movimento.

\subsection{Morte}

Molti dei personaggi non resteranno vivi sino alla fine della partita. Un personaggio può morire in diversi modi: condannato al rogo, sbranato dai Lupi, oppure ucciso da un'abilità o potere di altro tipo.

I personaggi morti possono comunque continuare a giocare e aiutare la propria fazione. Se una fazione vince, tutti i componenti di quella fazione vincono, indipendentemente dal fatto che siano vivi o morti.

Un personaggio morto non può votare, né usare la propria eventuale abilità; se ha un potere, può invece utilizzarlo. Può anche interagire con gli altri personaggi in tutti i modi non vietati dal regolamento, per esempio rivelando informazioni in suo possesso. Può inoltre essere riportato in vita dal Messia o divenire uno Spettro, se si verificano le opportune condizioni.

% Solo i Villani possono avere un potere, come spiegato nella Sezione~\ref{incantesimi}.

\subsection{Composizione del villaggio}

La composizione del villaggio è scelta dai GM e non è nota ai partecipanti. I ruoli sono assegnati casualmente ai giocatori.

I GM sono liberi, all'inizio della partita, di dare alcune indicazioni sulla composizione. Ad esempio, potrebbero dire: ``Su $40$ giocatori presenti, i Lupi sono più di $4$, i Villani non sono più di $11$ e ci sono almeno un Veggente ed un Espansivo''. Ciascuna di queste indicazioni può essere rivelata pubblicamente, oppure soltanto ad alcuni giocatori.

\subsection{Ruoli}
\label{ruoli}

I seguenti sono i ruoli che \emph{possono} comparire nel gioco, divisi per fazione. Per ogni ruolo sono specificati la fazione di appartenenza, il colore dell'aura (che può essere bianco o nero), la presenza o meno dell'attributo ``mistico'', ed eventuali tag.

\subsection*{Fazione dei Villici}

\begin{itemize}

      \item \textbf{Apprendista} [avvio, speciale]

            \emph{Cacciatore, Mago, Veggente, all'occorrenza anche Sequestratore, sai fare tantissime cose, nessuna delle quali particolarmente bene. Quando AlmaLaurea ti ha proposto una internship come "Assistant Werewolf Hunter" presso Casa del Diavolo (Perugia) a ben 400 Euro lordi al mese, non hai potuto dire no. Certo, pensavi fosse per un videogioco, e invece...}

            \begin{multicols}{3}
                  \begin{itemize}
                        \item \textbf{Aura:} Bianca.
                        \item \textbf{Frequenza:} 1/notte.
                        \item \textbf{Mistico:} No.
                  \end{itemize}
            \end{multicols}

            All'inizio della partita, l'Apprendista ottiene un elenco di abilità (eventualmente ripetute). Questo elenco non può comprendere le abilità di Lupo o Negromante.

            Ogni notte, l'Apprendista può selezionare un'abilità dal proprio elenco ed utilizzarla, specificando il personaggio sul quale la utilizza ed eventuali ulteriori informazioni. Se tale abilità viene utilizzata con successo, viene rimossa dall'elenco.

      \item \textbf{Bracconiere} [morte]

            \emph{TODO.}

            \begin{multicols}{3}
                  \begin{itemize}
                        \item \textbf{Aura:} Nera.
                        \item \textbf{Frequenza:} 1/2 notti.
                        \item \textbf{Mistico:} No.
                  \end{itemize}
            \end{multicols}

            Ogni due notti, il Bracconiere può scegliere un personaggio vivo e piazzare una trappola davanti alla sua casa. Per quella notte, se un Lupo utilizza la propria abilità su tale personaggio, il Lupo muore.

      \item \textbf{Cacciatore} [morte]

            \emph{Cerbiatti, coniglietti, allodole ed altri teneri animali del sottobosco per te non sono altro che selvaggina, cappotti, girocolli e scaldamani. E qualche volta un bel trofeo. Ti sei rifugiato in questo piccolo villaggio per sfuggire ad un gruppo di animalisti inferociti che aveva iniziato a pattugliare il tuo terreno di caccia preferito: la riserva naturale per la protezione del panda gigante.}

            \begin{multicols}{3}
                  \begin{itemize}
                        \item \textbf{Aura:} Bianca.
                        \item \textbf{Frequenza:} 1/partita.
                        \item \textbf{Mistico:} No.
                  \end{itemize}
            \end{multicols}

            Ogni notte, il Cacciatore può scegliere un personaggio vivo e sparargli. Tale personaggio muore.

            Dopo che utilizza la sua abilità con successo, il Cacciatore non può più utilizzare abilità per il resto della partita.

      \item \textbf{Contadino}

            \emph{Sei un contadino, ovvero un buzzurro, villano, bifolco, villico, bovaro, burino, eccetera. Hai essenzialmente due modalità di azione: ``influenzabile sempliciotto'', in cui credi supinamente a qualunque cosa ti dicano, e ``folla inferocita'', in cui brandisci un forcone e parti al linciaggio di chiunque ti capiti a tiro.}

            \begin{multicols}{3}
                  \begin{itemize}
                        \item \textbf{Aura:} Bianca.
                        \item \textbf{Frequenza:} n/a.
                        \item \textbf{Mistico:} No.
                  \end{itemize}
            \end{multicols}

            Il Contadino non ha alcuna abilità.

      \item \textbf{Divinatore} [avvio, info]

            \emph{Hai sempre pensato di essere un tipo qualunque, il contadino medio pronto a linciare alla cieca senza farsi troppe domande. Ma da quando nel villaggio hanno cominciato a succedere fatti strani continui a sognare tuo nonno che invece dei numeri del lotto ripete le solite tre frasi. Sei certo che almeno una sia vera e almeno una falsa. Ma riuscirai a scoprire in quali frasi si cela la verità?}

            \begin{multicols}{3}
                  \begin{itemize}
                        \item \textbf{Aura:} Bianca.
                        \item \textbf{Frequenza:} 1/2 notti.
                        \item \textbf{Mistico:} Sì.
                  \end{itemize}
            \end{multicols}

            All'inizio della partita, il Divinatore riceve un elenco di tre frasi nella forma ``Il personaggio X ha il ruolo Y''. Di queste tre frasi almeno una è vera ed almeno una è falsa.

            Ogni due notti, il Divinatore può scegliere un personaggio vivo ed un ruolo, ed effettuare una divinazione. Scopre se il personaggio scelto ha tale ruolo.

      \item \textbf{Esorcista} [blocco]

            \emph{L'apparizione della tua torreggiante figura avvolta in una tonaca nera inquieta tutti i passanti, tranne quando sei coperto di vomito verde. Negli ultimi anni hai fatto un giro panoramico del Lato Oscuro, alleandoti con lupi mannari, streghe, diavoli e perfino avvocati. Questo però non autorizza dei negromanti emopiagnoni a invadere il tuo territorio. Armato di bibbia rinforzata in acciaio e pistola ad acqua santa, sei pronto ad affrontare le armate delle tenebre. C'è un nuovo esorcista in città.}

            \begin{multicols}{3}
                  \begin{itemize}
                        \item \textbf{Aura:} Bianca.
                        \item \textbf{Frequenza:} 1/notte.
                        \item \textbf{Mistico:} Sì.
                  \end{itemize}
            \end{multicols}

            Ogni notte, l'Esorcista può scegliere un personaggio, vivo o morto, e benedire la sua casa. Per quella notte i poteri utilizzati su tale personaggio non hanno effetto. Questa abilità non è influenzata da poteri.

            L'Esorcista scopre se ha annullato l'effetto di almeno un potere.

      \item \textbf{Espansivo} [info]

            \emph{La tua voglia di amare gli altri è tanto potente che non puoi fare a meno di andare in giro ad abbracciare la gente. Ti presenti di notte nelle case altrui con una fetta di torta: qualcuno ti invita ad entrare e ti offre tè e biscotti, qualcun altro ti sbatte la porta in faccia. È ammirevole che tu non abbia ancora desistito, dopo aver cercato di abbracciare, nell'ordine, un lupo, un serial killer con l'ulcera, e Zannier.}

            \begin{multicols}{3}
                  \begin{itemize}
                        \item \textbf{Aura:} Bianca.
                        \item \textbf{Frequenza:} 1/2 notti.
                        \item \textbf{Mistico:} No.
                  \end{itemize}
            \end{multicols}

            Ogni due notti, l'Espansivo può scegliere un personaggio vivo e andare a trovarlo. Se tale personaggio ha aura bianca, questi scopre l'identità dell'Espansivo; altrimenti, l'abilità dell'Espansivo non ha effetto.

      \item \textbf{Guardia} [blocco, movimento]

            \emph{Provvisto di divisa nera, occhiali da sole a specchio, deltoidi quadruplici e cervello delle dimensioni di una nocciolina, non ti chiamano più gorilla da quando il povero primate, sentendosi offeso, ha protestato.}

            \begin{multicols}{3}
                  \begin{itemize}
                        \item \textbf{Aura:} Bianca.
                        \item \textbf{Frequenza:} 1/notte.
                        \item \textbf{Mistico:} No.
                  \end{itemize}
            \end{multicols}

            Ogni notte, la Guardia può scegliere un personaggio vivo e proteggerlo. La Guardia scopre tutti i movimenti generati dall'uso di abilità con il tag [male] sul personaggio scelto; le abilità che hanno generato tali movimenti non hanno effetto.

      \item \textbf{Mago} [info]

            \emph{Hai imparato i tuoi incantesimi da un manuale di D\&D. L'ultima volta che hai provato a fare una magia vera ti sei dato fuoco e hanno dovuto spegnerti con gli idranti. In compenso riesci perfettamente a tirare fuori un coniglio dal cappello, indovinare la prima carta di un mazzo con 53 carte uguali e segare in due una donna (ma non riattaccarla).}

            \begin{multicols}{3}
                  \begin{itemize}
                        \item \textbf{Aura:} Bianca.
                        \item \textbf{Frequenza:} 1/notte.
                        \item \textbf{Mistico:} Sì.
                  \end{itemize}
            \end{multicols}

            Ogni notte, il Mago può scegliere un personaggio, vivo o morto, e percepirne le capacità magiche. Scopre se tale personaggio è mistico.

      \item \textbf{Medium} [blocco, info]

            \emph{Hai costruito una carriera parlando con i morti, facendo soldi facili alle spalle dei boccaloni disposti a sganciare un bel po' di grano per sapere che la bisnonna Gertrude ha trovato la pace nell'aldilà. Andava tutto benissimo, almeno fino a quando i morti non hanno effettivamente cominciato a risponderti. Adesso che dici la verità, hai perso tutta la tua clientela: la nipote della povera Gertrude non voleva davvero sapere qual era l'ingrediente segreto del polpettone.}

            \begin{multicols}{3}
                  \begin{itemize}
                        \item \textbf{Aura:} Nera.
                        \item \textbf{Frequenza:} 1/2 notti.
                        \item \textbf{Mistico:} Sì.
                  \end{itemize}
            \end{multicols}

            Ogni due notti, il Medium può scegliere un personaggio morto e contattarne lo spirito. Scopre il ruolo di tale personaggio. Se tale personaggio utilizza un potere quella notte, tale potere non ha effetto.

      \item \textbf{Messia}

            \emph{Sei l'Eccelso, il Predestinato, il Santissimo (se sei religioso), l'Eletto (se sei molto nerd), l'Unto del Signore (lasciamo perdere che è meglio). Di fronte a te la luce avanza, le tenebre arretrano e riesci a far risollevare dalle tombe più morti che il regista medio di film di zombie.}

            \begin{multicols}{3}
                  \begin{itemize}
                        \item \textbf{Aura:} Bianca.
                        \item \textbf{Frequenza:} 1/partita.
                        \item \textbf{Mistico:} Sì.
                  \end{itemize}
            \end{multicols}

            Ogni notte, il Messia può scegliere un personaggio morto e resuscitarlo. Tale personaggio ritornerà in vita all'inizio del giorno seguente. Se il personaggio scelto è uno Spettro, l'abilità del Messia non ha effetto.

            Dopo che utilizza la sua abilità con successo, il Messia non può più utilizzare abilità per il resto della partita.

      \item \textbf{Stalker} [movimento]

            \emph{Più appiccicoso della melassa, più insistente di una canzone degli One Direction, più fastidioso di una colonia di formiche rosse insediata nelle mutande. Una volta puntata una vittima la segui ovunque, e questo include anche il gabinetto, le atroci riunioni del club del libro e i ricevimenti del professor Beltram.}

            \begin{multicols}{3}
                  \begin{itemize}
                        \item \textbf{Aura:} Bianca.
                        \item \textbf{Frequenza:} 1/2 notti.
                        \item \textbf{Mistico:} No.
                  \end{itemize}
            \end{multicols}

            Ogni due notti, lo Stalker può scegliere un personaggio vivo e pedinarlo. Scopre un eventuale movimento generato da tale personaggio.


      \item \textbf{Trasformista} [info]

            \emph{Anche il tuo numero di scarpe è un mistero. Nella vita hai cambiato più facce che cappelli e nemmeno tua madre potrebbe riconoscerti ormai, ammesso che tu ti ricordi chi è tua madre. Forse sei un agente segreto, un supereroe, o forse cerchi solo di rifarti una vita dove nessuno conosce i tuoi incresciosi trascorsi da necrofilo. La passione per i cimiteri però ti è rimasta e nel tempo libero ami passeggiare tra le tombe in cerca di una nuova identità.}

            \begin{multicols}{3}
                  \begin{itemize}
                        \item \textbf{Aura:} Bianca.
                        \item \textbf{Frequenza:} 1/notte.
                        \item \textbf{Mistico:} No.
                  \end{itemize}
            \end{multicols}

            Ogni notte, il Trasformista può scegliere un personaggio morto e rubarne l'identità. Se il personaggio scelto ha un ruolo appartenente alla Fazione dei Villici, il Trasformista lo scopre ed ottiene tale ruolo (includendo aura e misticità, e perdendo definitivamente ruolo e abilità precedenti); altrimenti, l'abilità del Trasformista non ha effetto.

      \item \textbf{Veggente} [info]

            \emph{Le tue mistiche abilità ti permettono di sapere chi è buono, chi è malvagio e chi non si è cambiato le mutande. Hai fatto i soldi evadendo regolarmente le tasse e facendo gli oroscopi a pagamento su Twitch. Proprio quando stavi per scappare in Polinesia, ecco che iniziano i morti squartati, gli spettri, i roghi e i linciaggi. Questo non l'avevi proprio previsto.}

            \begin{multicols}{3}
                  \begin{itemize}
                        \item \textbf{Aura:} Bianca.
                        \item \textbf{Frequenza:} 1/notte.
                        \item \textbf{Mistico:} Sì.
                  \end{itemize}
            \end{multicols}

            Ogni notte, il Veggente può scegliere un personaggio vivo e scrutarlo nella sua sfera di cristallo. Scopre il colore dell'aura di tale personaggio.

      \item \textbf{Voyeur} [movimento]

            \emph{Nascosto nel tuo cespuglio di fiducia con un binocolo dotato di sensori a infrarossi, ti diverti a spiare tutte le attività della notte: coppiette che copulano, riti satanici, lupi mannari che straziano orrendamente i corpi delle loro vittime, e ogni tanto un ubriaco che arriva barcollando e piscia nel tuo cespuglio.}

            \begin{multicols}{3}
                  \begin{itemize}
                        \item \textbf{Aura:} Bianca.
                        \item \textbf{Frequenza:} 1/2 notti.
                        \item \textbf{Mistico:} No.
                  \end{itemize}
            \end{multicols}

            Ogni due notti, il Voyeur può scegliere un personaggio vivo e spiarlo. Scopre tutti i movimenti diretti verso la casa di tale personaggio.

\end{itemize}


\subsection*{Fazione dei Villani}

All'inizio della partita, la Fazione dei Villani ottiene un Grimorio, ossia un elenco di Incantesimi (eventualmente ripetuti), chiamato Grimorio principale.

\begin{itemize}

      \item \textbf{Lupo} [male, morte]

            \emph{Di giorno solo le folte basette, le profonde occhiaie e la tendenza a grattarti col piede dietro l'orecchio tradiscono quello che diventi di notte: un terribile lupo mannaro! Sei una bestia assetata di sangue, nata per uccidere, e il tuo unico cruccio è che qualche volta i contadini armati di forconi ci azzeccano.}

            \begin{multicols}{3}
                  \begin{itemize}
                        \item \textbf{Aura:} Nera.
                        \item \textbf{Frequenza:} 1/notte.
                        \item \textbf{Mistico:} No.
                  \end{itemize}
            \end{multicols}

            Ogni notte, ciascun Lupo può scegliere un personaggio vivo e sbranarlo. Se tutti i Lupi che decidono di usare la propria abilità scelgono lo stesso personaggio, questi muore. Se almeno due Lupi indicano personaggi diversi, la loro abilità non ha effetto.

      \item \textbf{Alcolista} []

            \emph{Dopo una sfolgorante carriera prima a capo del governo lussemburghese, poi della Commissione Europea, sei stato cacciato via e rimpiazzato da una baciapantofole della Merkel, circondato da indifferenza ed impopolarità. Trasferitoti nel villaggio, i lupi mannari hanno ti hanno convinto a sfogare la tua rabbia verso il mondo aiutandoli nel loro piano di sterminio. Di giorno sfrutti la tua esperienza nelle trattative durante le assemblee per abbindolare i villici e trarli in inganno, mentre la notte vai in giro, annegando la disperazione nell'alcol.}

            \begin{multicols}{3}
                  \begin{itemize}
                        \item \textbf{Aura:} Bianca.
                        \item \textbf{Frequenza:} 1/notte.
                        \item \textbf{Mistico:} No.
                  \end{itemize}
            \end{multicols}

            Ogni notte, l'Alcolista può scegliere un personaggio, vivo o morto, e ubriacarsi sotto casa sua. L'abilità dell'Alcolista non ha effetto.

      \item \textbf{Assassino} [male, morte]

            \emph{John Kennedy, Paolo Borsellino e Kenny di South Park sono solo alcuni degli omicidi con cui non hai assolutamente niente a che fare, ma a cui ti ispiri per diventare famoso. Hai deciso di appendere il fucile al chiodo e ritirarti in campagna, ma la recente striscia di omicidi ti ha spinto a riprendere il tuo vecchio hobby.}

            \begin{multicols}{3}
                  \begin{itemize}
                        \item \textbf{Aura:} Nera.
                        \item \textbf{Frequenza:} 1/2 notti.
                        \item \textbf{Mistico:} No.
                  \end{itemize}
            \end{multicols}

            Ogni notte, l'Assassino può scegliere un personaggio vivo e un secondo personaggio, vivo o morto, e tendere un imboscata. Se il primo personaggio effettua un movimento verso il secondo, il primo personaggio muore.

      \item \textbf{Avvocato} []

            \emph{La tua specialità è parlare per tre ore senza mai riprendere fiato e senza dire assolutamente nulla. In un villaggio di contadini che non hanno mai sentito pronunciare una parola più lunga di tre sillabe, questo ti rende pressoché onnipotente.}

            \begin{multicols}{3}
                  \begin{itemize}
                        \item \textbf{Aura:} Nera.
                        \item \textbf{Frequenza:} 1/2 notti.
                        \item \textbf{Mistico:} No.
                  \end{itemize}
            \end{multicols}

            Ogni due notti, l'Avvocato può scegliere un personaggio vivo e far approvare per esso una legge \textit{ad personam}. Se durante il giorno successivo tale personaggio viene condannato al rogo, questi non muore.

      \item \textbf{Diavolo} [info]

            \emph{Tu sei l'Angelo Caduto, la Bestia, il Maligno! Il tuo numero fortunato è 666. Sei dotato di molteplici teste cornute e di tanto in tanto ti trasformi in drago, serpente, donna nuda o cantante rock. I tuoi unici problemi sono che la tua farina è andata di nuovo tutta in crusca e hai a casa un intero stock di pentole senza coperchi, che non sai come piazzare; per fortuna ne sai una più del diavolo, e questo significa che ne sai davvero tante.}

            \begin{multicols}{3}
                  \begin{itemize}
                        \item \textbf{Aura:} Nera.
                        \item \textbf{Frequenza:} 1/notte.
                        \item \textbf{Mistico:} Sì.
                  \end{itemize}
            \end{multicols}

            Ogni notte, il Diavolo può scegliere un personaggio vivo e un sottoinsieme dei ruoli, e scrutare fra le fiamme. Scopre se il ruolo del personaggio scelto appartiene a tale sottoinsieme.

      \item \textbf{Fantasma} [avvio]

            \emph{Anche se cammini ancora in questa valle di lacrime, ti senti molto portato per l'aldilà. Ti piace passeggiare tra cimiteri abbandonati, chiese sconsacrate e antiche rovine, ululando lugubremente nel vento. Se trovi delle catene, non puoi fare a meno di scuoterle. Non vedi l'ora di tirare le cuoia per far vedere a contadini e lupi mannari di cosa sei capace.}

            \begin{multicols}{3}
                  \begin{itemize}
                        \item \textbf{Aura:} bianca.
                        \item \textbf{Frequenza:} n/a
                        \item \textbf{Mistico:} Sì.
                  \end{itemize}
            \end{multicols}

            Il Fantasma non ha alcuna abilità. All'inizio della partita, il Fantasma ottiene un elenco di poteri (eventualmente ripetuti).

            Ogni notte, se il Fantasma è morto, può selezionare un potere dal proprio elenco ed utilizzarlo, specificando il personaggio sul quale lo utilizza ed eventuali ulteriori informazioni. Se tale potere viene utilizzato con successo, viene rimosso dall'elenco.

      \item \textbf{Fattucchiera} []

            \emph{Passi il sabato sera insieme alle tue colleghe, a far ribollire lingue di varano, peli di pipistrello e pollici di marinaio per scatenare letali epidemie di singhiozzo nel villaggio. Sei una crudele creatura del Maligno, ma la tua credibilità è sminuita dal fatto di chiamarti fata Zurlina. E il fatto che la gente scoppi a ridere al sentire il tuo nome, non fa che aumentare il tuo desiderio di trasformarli tutti in rospi.}

            \begin{multicols}{3}
                  \begin{itemize}
                        \item \textbf{Aura:} Nera.
                        \item \textbf{Frequenza:} 1/notte.
                        \item \textbf{Mistico:} Sì.
                  \end{itemize}
            \end{multicols}

            Ogni notte, la Fattucchiera può scegliere un personaggio, vivo o morto, ed un ruolo, e colpire tale personaggio con una fattura. Le abilità con il tag [info] utilizzate durante la notte ricevono responsi come se tale personaggio avesse il ruolo scelto.

      \item \textbf{Negromante} [avvio]

            \emph{Candele nere, pennarello indelebile rosso per disegnare pentagrammi, un chilo e mezzo di gioielleria etnica: sei pronto a risvegliare i morti! Adesso devi solo creare l'evento su Facebook e invitare tutti i tuoi amici occultisti al cimitero per una seduta spiritica con finale a sorpresa. Certo, i bigotti del villaggio non capiranno la tua particolare spiritualità, ma potrai sempre evocare nuovi spettri per dare una spinta alla tua vita sociale.}

            \begin{multicols}{3}
                  \begin{itemize}
                        \item \textbf{Aura:} Nera.
                        \item \textbf{Frequenza:} 1/notte.
                        \item \textbf{Mistico:} Sì.
                  \end{itemize}
            \end{multicols}

            All'inizio della partita, il Negromante ottiene un Grimorio.

            Ogni notte, il Negromante può scegliere un personaggio morto ed un Incantesimo scelto dal proprio Grimorio o dal Grimorio Maestro.

            Se il personaggio scelto è uno Spettro, l'Incantesimo selezionato viene attivato sullo Spettro scelto. Un eventuale Incantesimo precedentemente attivo su tale Spettro viene sostituito dall'Incantesimo selezionato.

            Se il personaggio scelto è un Villico, questi diventa uno Spettro, l'Incantesimo selezionato viene attivato su di esso, e comincia a giocare per la Fazione dei Villani. Dopo che utilizza questa abilità con successo, il Negromante non può più utilizzare abilità per il resto della partita.

            In tutti gli altri casi, l'abilità del Negromante non ha effetto.

            Se l'abilità del Negromante viene utilizzata con successo, l'Incantesimo selezionato viene rimosso dal Grimorio.


      \item \textbf{Sequestratore} [blocco, male]

            \emph{Sei abitualmente domiciliato in una grotta nei pressi di Sussubbiarrutarruggiu (Sardegna), da dove riemergi per rapire facoltosi turisti in cerca di un cospicuo riscatto. Ma poi li rilasci in cambio di un biglietto del Lotto che ti è stato assicurato essere vincente, il che dimostra quanto vivere in una grotta non giovi alla salute mentale della gente.}

            \begin{multicols}{3}
                  \begin{itemize}
                        \item \textbf{Aura:} Nera.
                        \item \textbf{Frequenza:} 1/notte.
                        \item \textbf{Mistico:} No.
                  \end{itemize}
            \end{multicols}

            Ogni notte, il Sequestratore può scegliere un personaggio vivo e rapirlo. Se tale personaggio effettua un movimento, il Sequestratore scopre tale movimento, e l'abilità che lo ha generato non ha effetto.

      \item \textbf{Stregone} [blocco]

            \emph{La tua passione per l'occulto, il tuo odio verso i bifolchi del villaggio, e la tua cotta per la fata Zurlina ti hanno spinto ad imparare la magia nera e diventare uno Stregone. Passi il tuo tempo a farti beffe dei contadini ignoranti con la tua magia, sperando così di fare colpo sulla tua bella.}

            \begin{multicols}{3}
                  \begin{itemize}
                        \item \textbf{Aura:} Nera.
                        \item \textbf{Frequenza:} 1/notte.
                        \item \textbf{Mistico:} Sì.
                  \end{itemize}
            \end{multicols}

            Ogni notte, lo Stregone può scegliere un personaggio, vivo o morto, e lanciare un sortilegio sulla sua casa. Per quella notte, se un qualsiasi altro personaggio utilizza la propria abilità o potere sul personaggio scelto dallo Stregone, tale abilità o potere non ha effetto. %Lo Stregone scopre tutti i movimenti diretti verso la casa di tale personaggio.
\end{itemize}


\subsubsection*{Spettri e Incantesimi}
\label{spettri}
\label{incantesimi}

Lo Spettro non è un ruolo, ma uno status che può o meno essere assegnato ad un personaggio morto. In particolare, all'inizio della partita non vi è alcuno Spettro.

Ci sono due modi in cui un personaggio morto può divenire uno Spettro.

\begin{itemize}
      \item Quando un Villico muore, c'è la possibilità che questi divenga automaticamente uno Spettro, in un modo che dipende solamente da quanti altri Villici sono morti fino a quel momento (in un modo segretamente deciso dai GM prima dell'inizio della partita). In questo caso, nessun Incantesimo è attivo su tale Spettro; i Negromanti potranno attivarne uno a partire dalla notte successiva.

      \item Un Villico morto può divenire uno Spettro successivamente, in seguito all'uso dell'abilità di un Negromante. In questo caso un Incantesimo viene immediatamente attivato su tale Spettro, come riportato nella descrizione del Negromante.
\end{itemize}

Nel momento in cui un personaggio diviene uno Spettro, mantiene il suo ruolo (includendo aura e misticità), diventa un Villano, gli viene comunicata l'identità degli altri Villani, e a questi viene comunicata la sua identità.

Quando un personaggio riceve lo status di Spettro, questo rimane morto, a meno che su di esso non venga contemporaneamente assegnato l'Incantesimo della Vita. In particolare uno Spettro morto non può essere ucciso, e il villaggio non riceve alcuna informazione sull'assegnazione dello status di Spettro.

Un Villico non può diventare uno Spettro durante la stessa notte in cui la Fazione dei Villani viene sconfitta.

Uno Spettro che non abbia un Incantesimo attivo ha un potere: ogni notte può selezionare un Incantesimo dal Grimorio principale, che viene attivato sullo Spettro. Tale Incantesimo viene rimosso dal Grimorio principale.

La maggior parte degli Incantesimi assegna un potere ad uno Spettro. Se non diversamente riportato, la descrizione dell'Incantesimo coincide con quella del potere ottenuto dallo Spettro.

\begin{itemize}

      \item {\bf Assoluzione}. Ogni due notti, lo Spettro può scegliere un personaggio vivo e infondere un'aura di innocenza intorno a lui. Durante il giorno successivo, se tale personaggio viene condannato al rogo, questi non muore; se ciò accade, l'Incantesimo dell'Assoluzione viene disattivato.

      \item {\bf Confusione}. Ogni notte, lo Spettro può scegliere un personaggio, vivo o morto, e generare un'aura di confusione attorno al personaggio scelto. Per quella notte, il colore dell'aura di tale personaggio viene invertito.

      \item {\bf Invisibilità} [movimento]. Ogni due notti, lo Spettro può scegliere un personaggio vivo e renderlo invisibile. Un eventuale movimento generato da tale personaggio viene eliminato. Eventuali abilità di tale personaggio non sono influenzate.

      \item {\bf Morte} [morte]. Ogni notte, lo Spettro può scegliere un personaggio vivo e un ruolo, e spaventarlo a morte. Se tale personaggio ha il ruolo selezionato, questi muore. L'Incantesimo della Morte viene disattivato.

      \item {\bf Occultamento} [blocco]. Ogni due notti, lo Spettro può scegliere un personaggio, vivo o morto, e creare attorno alla sua casa una nebbia magica. Per quella notte, se un qualsiasi personaggio utilizza la propria abilità sul personaggio scelto dallo Spettro, tale abilità non ha effetto.

      \item {\bf Telepatia} [info]. Ogni due notti, lo Spettro può scegliere un personaggio vivo e un ruolo, e leggerne la mente. Scopre se il personaggio scelto ha tale ruolo; in caso di responso positivo, ottiene tutte le informazioni che tale personaggio riceve durante l'alba successiva.

      \item {\bf Visione} [info]. Ogni notte, lo Spettro può scegliere un personaggio, vivo o morto, e spiarne la casa. Scopre l'elenco di tutte le abilità (eventualmente ripetute) che sono state utilizzate quella notte su quel personaggio (ma non scopre l'identità dei giocatori che le hanno utilizzate).

      \item {\bf Vita}. Lo Spettro non ottiene alcun potere. Durante l'alba successiva, lo Spettro ritorna in vita mantenendo ruolo ed abilità, ma continuando a giocare con i Villani. L'Incantesimo della Vita viene disattivato se si verifica una condizione per cui tale personaggio dovrebbe morire (per esempio perché bruciato sul rogo). Se l'Incantesimo della Vita viene disattivato, lo Spettro muore.

\end{itemize}

Quando un Incantesimo viene attivato o disattivato su di uno Spettro, questi riceve una notifica.

\subsection{Utilizzo di abilità, poteri, e fallimento}
\label{fallimento}

\paragraph{Utilizzo di abilità e poteri}

Se un personaggio con un'abilità o un potere attivabile ogni due notti (oppure una volta a partita) lo utilizza con successo, non potrà utilizzare alcuna abilità o potere nella notte successiva (o nel resto della partita).

I personaggi che ottengono informazioni di qualsiasi tipo su altri personaggi, le ottengono riguardo alla condizione in cui questi si trovavano al termine del giorno precedente, anche se tali informazioni cambiano nel corso della notte.
Ad esempio, se un Trasformista viene visitato da un Veggente durante la notte in cui utilizza la propria abilità con successo su un Medium, il Veggente lo vedrà comunque con aura bianca, perché la trasformazione avviene al termine della notte.

\paragraph{Fallimento} Ci sono alcuni modi di utilizzare la propria abilità o il proprio potere che sono a priori non consentiti, e pertanto non risultano nemmeno tra le scelte disponibili sull'interfaccia web. Ad esempio, il Medium non può scegliere di usare la propria abilità su un personaggio vivo; un Messia non può provare ad attivare la propria abilità se l'ha già utilizzata con successo in una notte precedente.

Anche escludendo i casi menzionati sopra, non sempre un personaggio riesce ad usare la propria abilità o potere con successo. Ci sono essenzialmente tre possibili ragioni per cui questo può accadere.

\begin{itemize}
      \item Blocco da parte di un altro personaggio. Ad esempio, si può essere rapiti da un Sequestratore, bloccati da un Esorcista o dallo Spettro su cui è attivo l'Incantesimo dell'Occultamento.
      \item Restrizioni all'utilizzo di abilità o potere. Ad esempio, un Lupo non riesce a uccidere se nella stessa notte un altro Lupo cerca di assassinare un personaggio diverso; l'abilità del Trasformista non può essere usata su un Lupo.
      \item Contraddizioni nell'utilizzo di abilità o poteri. Può succedere che non ci sia alcun modo logico di risolvere le azioni dei personaggi nel corso di una notte (vedi le FAQ per una discussione più approfondita). Questo accade per esempio se tre Sequestratori si sequestrano in modo ciclico.
\end{itemize}

Nel caso in cui un personaggio cerchi di utilizzare la propria abilità o potere ma l'utilizzo non ha effetto, il personaggio riceve una notifica di fallimento. Tale notifica non precisa la ragione del fallimento. Un utilizzo a vuoto (ad esempio un Cacciatore che utilizzi la propria abilità su un personaggio diverso da quello scelto dai Lupi, oppure una Fattucchiera che assegna l'apparenza di Veggente ad un Veggente) non costituisce un fallimento.

Un personaggio vivo che cerchi di utilizzare la propria abilità genera un movimento anche se questa non ha effetto.

\section{Credits}
\label{credits}

Si ringraziano queste persone per aver contribuito in vari ruoli a questa edizione di \emph{Lupus in tempo reale} o alle precedenti, in ordine circa cronologico.

\subsection*{Regolamento edizione corrente}

Claudio Afeltra, Lorenzo Bresolin, Mariastella Cascone, Alessandro Iraci, Giovanni Italiano, Leonardo Tolomeo.

% \subsection*{Consigli di gioco} %TODO

\subsection*{Regolamento edizioni passate}

Luca Ghidelli, Alessandro Iraci, Giovanni Mascellani, Giovanni Paolini, Leonardo Tolomeo, Matteo Migliorini, Giovanni Italiano, Claudio Afeltra, Edoardo Rizzi, Giovanni Interdonato, Lorenzo Pierro.

\subsection*{Server di gioco}

\textbf{Struttura principale del server:} Giovanni Mascellani, Giovanni Paolini, Matteo Migliorini.

\textbf{Altri contributi:} Flavio Ascari, Manuele Cusumano, Alessandro Iraci, Roberto Tocco.

\subsection*{Concept e prime edizioni}

Andrea Caleo, Francesco Guatieri.

\end{document}
